% !TEX root = main.tex
Il segnale $y(t)$ assegnato è un segnale di potenza. La potenza media $P_y$ è calcolabile nel dominio del tempo tramite il suo integrale quadratico su un periodo $T_0$.

Normalizzando l'asse con $x = \frac{2t}{T_0}$ (con $dt = \frac{T_0}{2}dx$), l'integrale risulta:
\[ P_y = 2 \int_{0}^{0.5} [3(1-x) + 1.5]^2\, dx + 2 \int_{0.5}^{1} [3(1-x)]^2\, dx \]
\[ P_y = \left[ -\frac{(4.5 - 3x)^3}{9} \right]_0^{0.5} + \left[ -\frac{(3 - 3x)^3}{9} \right]_{0.5}^1 = \frac{64.125}{9} + \frac{3.375}{9} = \textbf{7.5 V}^2 \]

La potenza media dissipata su un carico unitario è di $\mathbf{7.5 \text{ W}}$.

\subsection{Occupazione in Banda al 99\%}
Definiamo la banda essenziale come la soglia di frequenza che racchiude il 99\% della potenza totale calcolata. Sommando la potenza delle singole righe armoniche $|Y_k|^2$:
\begin{itemize}
    \item \textbf{DC ($k=0$):} $P_{DC} = (2.25)^2 = 5.0625 \text{ V}^2$
    \item \textbf{Frequenza $f_1$ (6 kHz):} $P_{1} = 2 \times (1.0854)^2 \approx 2.356 \text{ V}^2$
    \item \textbf{Frequenza $f_3$ (18 kHz):} $P_{3} = 2 \times (0.091)^2 \approx 0.016 \text{ V}^2$
\end{itemize}

Verifichiamo la potenza cumulativa:
\[ \frac{P_{DC} + P_1}{P_{tot}} = \frac{7.418}{7.5} = 98.9\% \]
\[ \frac{P_{DC} + P_1 + P_3}{P_{tot}} = \frac{7.434}{7.5} = 99.1\% \]

La soglia del 99\% viene superata alla terza armonica, identificando una banda di occupazione pari a $\boxed{B_{99\%} = 18 \text{ kHz}}$.
